\section{Causal evidence for organizational benefit}\label{sec:C24:benefit}
Organizational benefit is not directly measured by model benchmark accuracy. An agent can produce useful drafts while increasing review effort, or reduce time on routine work while creating rare costly errors. Deployment studies must define the outcome and the accounting boundary. Candidate outcomes include verified task completion, time to accepted artifact, reviewer effort, rework, error severity, and downstream process effects.

The Knowledge Activation paper proposes structured institutional knowledge units and reports an enterprise case study~\cite{corpus260314805}. Its developer survey supports claims about perceived experience in that setting. Self-reported time savings are not instrumented productivity measurements, and concurrent model or tooling changes complicate causal attribution. The useful architectural contribution can therefore be discussed without promoting the case study into a universal magnitude claim.

A stronger study randomizes eligible work units or uses a carefully justified staggered rollout. The unit of assignment should minimize interference: engineers sharing skills and repositories can transfer treatment effects to nominal controls. Measure baseline workload, task difficulty, and prior tool exposure. Record adoption and noncompliance, because assigning an agent does not guarantee that workers use it. An intention-to-treat analysis and a usage-conditioned analysis answer different questions and should not be conflated.

Learning effects require longitudinal observation. Early users may spend extra time configuring context and tools; later users may benefit from accumulated skills. Conversely, initial novelty can inflate perceived benefit while maintenance costs appear later. Report trajectories rather than only a single post-deployment mean. Include unsuccessful and abandoned tasks, since excluding them can make a system appear faster by conditioning on completion.

External validity depends on organizational knowledge quality. A system grounded in accurate, maintained procedures may perform differently from one grounded in contradictory documentation. The treatment therefore includes knowledge engineering and governance, not merely an agent interface. An ablation can hold the model and harness fixed while varying institutional context quality. This identifies whether observed gains come from the policy, better information, or their interaction.

